\documentclass[12pt,aspectratio=169]{beamer}
\input{../../mybeamer}

\title{Classes 5 \& 6: Capacitors}
\subtitle{AP Physics C: Electricity and Magnetism}
\input{../../me}
\input{../../mycommands}


\begin{document}

\begin{frame}
  \maketitle
\end{frame}


\section{Parallel-Plate Capacitors}

\begin{frame}{Electric Field and Electric Potential Difference}
  Recall that the relationship between electrostatic force ($\vec F_q$) and
  electric potential energy ($U_q$) can be expressed using definition of
  mechanical work and the fundamental theorem of calculus:

  \eq{-.1in}{
    \Delta U_q=-\int\vec F_q\cdot\dl\vec r\quad\quad
    \vec F_q=-\nabla U_q=-\diffp{U_q}r\hat r
  }

  Dividing both sides of the equations by $q$, we get the relationship between
  electric field ($\vec E$), electric potential ($V$) and electric potential
  difference ($\Delta V$):
  
  \eq{-.1in}{
    \Delta V=-\int\vec E\cdot\dl\vec r\quad\quad
    \boxed{\vec E=-\nabla V=-\diffp Vr\hat r}
  }

  This relationship holds regardless of the charge configuration.
\end{frame}



\begin{frame}{Electric Field and Electric Potential Difference}
  \begin{center}
    \begin{tikzpicture}[thick]%[xscale=.7,yscale=.9]
      \foreach \x in {.4,.8,1.2,...,7.6}{
        \draw[vectors,red] (\x,1.4)--(\x,.7);
        \draw[very thick,red] (\x,1.4)--(\x,.2);
      }
      \foreach\x in {.4,.8,...,7.9}{ %{.8,1.6,2.4,...,7.2}{
        \node at (\x,1.5) {\small$\bm +$};
        \node at (\x,.1) {\small$\bm -$};
      }
      
      \draw[vectors,red] (0,1.5) ..controls(-.15,1.2).. (-.2,.7);
      \draw[very thick,red] (-.2,.8) ..controls(-.15,.4).. (0,.1);

      \draw[vectors,red] (8,1.5) ..controls(8.15,1.2).. (8.2,.7);
      \draw[very thick,red] (8.2,.8) ..controls(8.15,.4).. (8,.1);

      \draw rectangle +(8,.2);
      \draw (0,1.4) rectangle +(8,.2);
    \end{tikzpicture}
  \end{center}
  Recall that, for two charged parallel plates (where the area $A$ of the plate
  is larger compared to plate separation $d$), the electric field $E$ between
  the plate is uniform. There are two ways to express the electric field
  strength:

  %, and the relationship between electric field and
  %potential difference simplifies to:

  \eq{-.1in}{
    \boxed{E=\frac{\Delta V}d}
    \quad\text{or}\quad
    \boxed{E=\frac{\sigma}{\varepsilon_0}}
    %\boxed{\Delta V=Ed}
  }
\end{frame}



\begin{frame}{Capacitors}
  \textbf{Capacitors} is a device that stores energy in an electric field. The
  simplest form of a capacitor is a set of closely spaced parallel plates:
  \begin{center}
    \pic{.5}{cap19}
  \end{center}
  When the plates are connected to a battery, the battery transfer charges to
  the plates until the voltage $V$ equals the battery terminals. After that,
  one plate has charge $+Q$; the other has $-Q$.
\end{frame}



\begin{frame}{Parallel-Plate Capacitors}
  As we have seen already, the (uniform) electric field between two parallel
  plates is proportional to the charge density $\sigma$, which is the charge
  $Q$ divided by the area of the plates $A$:

  \eq{-.1in}{
    E=\frac{\color{red}\sigma}{\varepsilon_0}=
    \frac{\color{red}Q}{{\color{red}A}\varepsilon_0}
  }
  
  Equating the two expressions for electric field, we find a linear
  relationship between the plate voltage $V$ and charges across the
  plates and the voltage:

  \eq{-.1in}{
    E=\frac Q{A\varepsilon_0}=\frac{\Delta V}d
    \quad\longrightarrow\quad
    \boxed{Q=\left[\frac{A\varepsilon_0}d\right]\Delta V}
  }
\end{frame}



\begin{frame}{Parallel-Plate Capacitors}
  Since area $A$, distance of separation $d$ and the vacuum permittivity
  $\varepsilon_0$ are all constants, the relationship between charge $Q$ and
  voltage $\Delta V$ is \emph{linear}. The constant is called the
  \textbf{capacitance} $C$, defined as:

  \eq{-.1in}{
    \boxed{C=\frac Q{\Delta V}}
  }

  For parallel plates:

  \eq{-.1in}{
    \boxed{C=\frac{A\varepsilon_0}d\quad\text{parallel plate}}
  }

  The unit for capacitance is a \textbf{farad} (named after Michael Faraday),
  where $\SI1\farad=\SI1{\coulomb\per\volt}$.
\end{frame}



\begin{frame}{Storage of Electrical Energy}
  When charging up a capacitor, imagine positive charges moving from the
  negatively charged plate to the positively charged plate
  \begin{center}
    \pic{.45}{slide14}
  \end{center}
  %Initially neither plates are charged, so moving the first charge takes very
  %little work; as the electric field builds, more and more work needs to be
  %done
\end{frame}



\begin{frame}{Storage of Electrical Energy}
  \begin{center}
    \pic{.35}{slide14}
  \end{center}
  In the beginning---when the plates aren't charged---moving an infinitesimal
  charge $\dl q$ across the plates, the infinitesimal work done $\dl U$ is very
  small and related to the capacitance by:

  \eq{-.1in}{
    \dl U=V\dl q=\frac qC\dl q
  }
  
  As the electric field begins to form between plates, more and more work
  is required to move the charges.
\end{frame}



\begin{frame}{Storage of Electrical Energy}
  To fully charge the plates, the total work $U_c$ is the integral:

  \eq{-.1in}{
    U_c=\int\dl U=\int_0^Q\frac qC\dl q=\frac12\frac{Q^2}C
  }

  The work done is stored as a potential energy inside the capacitor. There are
  different ways to express $U_c$ using definition of capacitance:

  \eq{-.1in}{
    \boxed{U_c=\frac12\frac{Q^2}C=\frac12QV=\frac12CV^2}
  }
\end{frame}



\section{Cylindrical Capacitors}

\begin{frame}{Cylindrical Capacitors}
  Not all capacitors are parallel plates. Cylindrical capacitors are also
  popular:
  \begin{columns}
    \column{.23\textwidth}
    \begin{tikzpicture}[scale=.6]
      \fill[cyan!30,opacity=.2] (-2,8)--(-2,0) arc (180:0:2.5 and 1)--(3,8)
      arc (0:180:2.5 and 1);
      \draw[cyan,thick] (-2,8)--(-2,0);
      \draw[cyan,thick] (3,0)--(3,8) arc (0:180:2.5 and 1);
      
      \fill[violet,opacity=.1] (-1,8)--(-1,0) arc (180:0:1.5 and .5) -- (2,8)
      arc (0:180:1.5 and .5) -- (-1,8);
      \draw[dash dot,thick,violet] (-1,8)--(-1,0) arc (180:0:1.5 and .5)--
      (2,8) arc (0:180:1.5 and .5);

      \fill[orange!50, opacity=.6] (0,8)--(0,0) arc (180:0:.5 and .2) -- (1,8)
      arc (0:180:.5 and .2) -- (0,8);
      \draw[orange,thick] (0,0)--(0,8) arc (180:0:.5 and .2)--(1,0);
      
      \fill[orange,opacity=.8] (0,8)--(0,0) arc (180:360:.5 and .2) -- (1,8)
      arc (0:-180:.5 and .2) -- (0,8);

      \fill[violet,opacity=.1] (-1,8)--(-1,0) arc (180:360:1.5 and .5)--(2,8)
      arc (0:-180:1.5 and .5) -- (-1,8);
      \draw[dash dot,thick,violet] (-1,8) arc (180:360:1.5 and .5);

      \fill[cyan!30] (-2,8)--(-2,0) arc (180:360:2.5 and 1)
      --(3,8) arc (0:-180:2.5 and 1);
      \draw[cyan] (-2,8)--(-2,0) arc (180:360:2.5 and 1)
      --(3,8) arc (0:-180:2.5 and 1);

      \draw[|<->|] (3.4,8)--(3.4,0) node[midway,fill=black!2]{$\ell$};
      \draw[axes,rotate around={10:(.5,8)}](.5,8)--(.95,8) node[right]{$a$};
      \draw[axes,rotate around={-30:(.5,8)}](.5,8)--(2.15,8) node[pos=1.2]{$b$};
    \end{tikzpicture}

    \column{.15\textwidth}
    \centering
    \begin{tikzpicture}[scale=.8]
      \draw[thick,cyan,fill=cyan!30] circle (1.5);
      \draw[thick,cyan,fill=black!2] circle (1.45);
      \foreach\theta in {0,30,...,359}
      \draw[vectors,red,rotate=\theta] (.5,0)--(1.45,0);
      \draw[thick,orange,fill=orange!50] circle (.5);
      \draw[axes,rotate=30] (0,0)--(.5,0) node[midway,below]{$a$};
      \draw[axes,rotate=110] (0,0)--(1.45,0) node[pos=1.2]{$b$};
      \node[red] at (-.2,-1){$\vec E$};
      %\draw[violet,thick,dash dot] circle (1);
      %\draw[thick,<-,violet] (-.707,-.707)--+(-1,-1)
      %node[below=-3]{Gaussian Surface};
    \end{tikzpicture}\\
    {\footnotesize End view\par}
    
    \column{.62\textwidth}
    \begin{itemize}
    \item The capacitor has length $\ell$ which is much larger than the radii
      of the inner \& outer cylinders ($a$, $b$)
    \item Inner cylinder has total charge $Q$
    \item Outer cylinder has total charge $-Q$
    \item Inside the capacitor, the electric field in the radial direction
    \item Outside of the capacitor, there is no electric field
    \end{itemize}
  \end{columns}
\end{frame}



\begin{frame}{Cylindrical Capacitors: Electric Field}
  \begin{columns}
    \column{.33\textwidth}
    \begin{tikzpicture}
      \draw[thick,cyan,fill=cyan!30] circle (1.5);
      \draw[thick,cyan,fill=black!2] circle (1.45);
      \foreach \theta in {0,30,...,359}
      \draw[vectors,red,rotate=\theta] (.5,0)--(1.45,0);
      \draw[thick,orange,fill=orange!50] circle (.5);
      \draw[axes,rotate=30] (0,0)--(.5,0) node[midway,below]{$a$};
      \draw[axes,rotate=110] (0,0)--(1.45,0) node[midway,right]{$b$};
      \node[red] at (-.2,-1){$\vec E$};
      \draw[violet,thick,dash dot] circle (1);
      \draw[thick,<-,violet] (-.707,-.707)--+(-1,-1)
      node[below=-3]{Gaussian Surface};
    \end{tikzpicture}
    %\pic{1.05}{cylindrical-capacitor}

    \column{.67\textwidth}
    We can use Gauss's law to find the electric field between the cylinders, by
    placing a Gaussian surface of radius $r$ between the cylinders: 

    \eq{-.1in}{
      \oint\vec E\cdot\dl\vec A=2\pi r\ell E=\frac Q{\varepsilon_0}
    }
    
    which gives the expression:

    \eq{-.1in}{
      E=\frac Q{2\pi r\ell\varepsilon_0}\quad\text{or}\quad
      E=\frac\lambda{2\pi r\varepsilon_0}
    }
    
    where $\lambda=Q/\ell$ is the linear charge density
  \end{columns}
\end{frame}




\begin{frame}{Cylindrical Capacitors: Voltage Across the Cylinders}
  Integrating the eclectic field to get voltage across the plates:

  \eq{-.1in}{
    |\Delta V|= \int_a^b E\dl r=\frac Q{2\pi\ell\varepsilon_0}\int_b^a\frac{\dl r}r
    =\frac Q{2\pi\ell\varepsilon_0}\ln\left[\frac ba\right]
  }

  Like parallel plates, the relationship between voltage and charge is still
  linear, but in this case, the capacitance is defined as:

  \eq{-.1in}{
    C=\frac Q{\Delta V}=\frac{2\pi\ell\varepsilon_0}{\ln(b/a)}
    \quad\to\quad
    \boxed{\frac C\ell=\frac{2\pi\varepsilon_0}{\ln(b/a)}}
  }

  The capacitance of a cylindrical capacitor is often expressed by $C/\ell$
  (with unit \si{\farad\per\metre}). Like the parallel-plate capacitor, the
  capacitance of the cylindrical capacitor also only depends on the geometry
  (i.e.\ the raidii $a$ and $b$) and the permittiviy.
\end{frame}



\section{Spherical Capacitor}

\begin{frame}{Spherical Capacitor}
  Two spherical charges can also be arranged to act as a capacitor. In the
  configuration shown below, the inner sphere carries a charge of $+Q$, while
  the outer spherical shell carries a charge of $-Q$.
  \begin{columns}[T]
    \column{.5\textwidth}
    \centering
    \begin{tikzpicture}
      \shade[ball color=cyan] circle (.7);
      \draw[thick,cyan] circle (.7);
      \shade[ball color=magenta,opacity=.3] circle (2);
      \draw[thick,magenta] circle (2);
      \draw[axes,rotate=100] (0,0)--(.7,0) node[midway,right]{$a$};
      \draw[axes,rotate=-130] (0,0)--(2,0) node[midway,below]{$b$};
    \end{tikzpicture}
    
    \column{.5\textwidth}
    \begin{tikzpicture}
      \draw[thick,magenta,fill=magenta!30] circle (2);
      \draw[thick,magenta,fill=black!2] circle (1.95);
      \draw[thick,cyan,fill=cyan!30] circle (.7);
      \draw[axes,rotate=100] (0,0)--(.7,0) node[midway,right]{$a$};
      \draw[axes,rotate=-120] (0,0)--(1.95,0) node[midway,below]{$b$};
      \foreach \theta in {15,45,...,360}{
        \draw[vectors,orange,rotate=\theta] (.7,0)--(1.95,0);
      }
      \draw[thick,dashed,red] circle (1);
      \draw[thick] (0,-2)--(0,-2.5)--(2.5,-2.5)
      to[battery,l_=$\mathcal E$] (2.5,0)--(2.05,0) arc (0:-180:.08)--(.7,0);
    \end{tikzpicture}
  \end{columns}

  We are interested in the electric field between the spheres.
\end{frame}



\begin{frame}{Spherical Capacitor}
  Whave already determined, using Gauss's law, the magnitude of the electric
  field between the inner and outer spheres:
  
  \eq{-.1in}{
    E=\frac1{4\pi\varepsilon_0}\frac Q{r^2}
  }

  Therefore the electric potential difference between the inner sphere ($r=a$)
  and the outer spherical shell $r=b$) is given by:

  \eq{-.1in}{
    |\Delta V|=\int_a^bE\dl r=\frac Q{4\pi\varepsilon_0}\int_a^b\frac{\dl r}{r^2}
    =\frac Q{4\pi\varepsilon_0}\left(\frac1a-\frac1b\right)
  }
\end{frame}



\begin{frame}{Capacitance}
  Regardless of the geometry of the capacitor, the electric field is always
  proportional to the charge, i.e.:

  \eq{-.2in}{ E\propto Q }

  and therefore the voltage will always be proportional to charge as well.
\end{frame}



\section{Practical Capacitors}

\begin{frame}{Practical Capacitors}
  \begin{columns}
    \column{.3\textwidth}
    \pic{1.15}{Figure_20_05_05a}

    \column{.7\textwidth}
    \begin{itemize}
    \item Capacitors (both parallel-plate and cylindrical) are very common in
      electric circuits, but the vacuum between the plates is not very effective
    \item Instead, a non-conducting \textbf{dielectric} material is inserted
      between the plates
    \item When the plates are charged, the electric field of the plates
      polarizes the dielectric.
    \item The polarization  produces an electric field that opposes the field
      from the plates, therefore reduces the effective voltage, and increasing
      the capacitance
    \end{itemize}
  \end{columns}
\end{frame}



\begin{frame}{Dielectric Constant}
  If electric field without dielectric is $E_0$, then $E$ in the dielectric is
  reduced by $\kappa$, the \textbf{dielectric constant}:

  \eq{-.1in}{
    \boxed{\kappa=\frac{E_0}E}
  }

  The capacitance of the plates with the dielectric is now amplified by the
  same factor $\kappa$:

  \eq{-.1in}{
    \boxed{C=\kappa C_0}
  }

  We can also view the dielectric as something that increases the
  \emph{effective permittivity}:
  
  \eq{-.1in}{
    \boxed{\varepsilon=\kappa\varepsilon_0}
  }
\end{frame}



\begin{frame}{Dielectric Constant}
  The dielectric constants of commonly used materials are:
  \begin{center}
    \begin{tabular}{l|l}
      \rowcolor{pink}
      \textbf{Material} & $\kappa$ \\ \hline
      Air         & \num{1.00059} \\
      Bakelite    & \num{4.9} \\
      Pyrex glass & \num{5.6} \\
      Neoprene    & \num{6.9} \\
      Plexiglas   & \num{3.4} \\
      Polystyrene & \num{2.55} \\
      Water (\SI{20}\celsius) & \num{80} 
    \end{tabular}
  \end{center}
\end{frame}



\begin{frame}{Notes About Storage of Electric Energy}
  The work done (i.e.\ the energy stored in the capacitor) is inversely
  proportional to the capacitance:

  \eq{-.1in}{
    \dl U=V\dl q=\frac qC\dl q
  }

  \begin{itemize}
  \item The presence of a dielectric \emph{increases} the capacitance; therefore
    the work (and potential energy stored) to move the charge $\dl q$
    \emph{decreases} with the dielectric constant $\kappa$
  \item After the capacitor is charged, removing the dielectric material from
    the capacitor plates will require additional work.
  \end{itemize}
\end{frame}



\section{Capacitors in Circuit}

\begin{frame}{Capacitors in Electric Circuits}
  Capacitors are an important part of an electric circuits because it stores
  energy in the electric field. Suppose we replace a battery with a capacitor
  in a basic resistive circuit:
  \begin{center}
    \begin{tikzpicture}
      \draw[thick] (0,0) to[C=$C$] (0,2)--(2,2) to[R=$R$] (2,0)--(0,0);
    \end{tikzpicture}
  \end{center}
  \begin{itemize}
  \item The capacitor acts like a voltage source.
  \item Unlike a battery, the voltage decreases as it drives a current, and the
    charge across the capacitor plates decreases.
  \end{itemize}
\end{frame}



\begin{frame}{Capacitors in Parallel}
  \begin{center}
    \begin{tikzpicture}[scale=1.2,thick]
      \draw (0,1) to[short,o-] (1,1) to[C=$C_1$] (1,0) to[short,-o] (0,0);
      \draw (1,1)--(3,1) to[C=$C_2$] (3,0)--(1,0);
      \draw (3,1)--(5,1) to[C=$C_3\ldots$] (5,0)--(3,0);
    \end{tikzpicture}
  \end{center}
  From the voltage law, we know that the voltage across all the capacitors are
  the same, i.e.\ $V_1=V_2=V_3=\cdots=V$. We can express the total charge
  $Q_\text{tot}$ stored across all the capacitors in terms of capacitance and
  this common voltage $V$: 

  \eq{-.1in}{
    Q_\text{tot}=Q_1+Q_2+Q_3+\cdots=C_1V+C_2V+C_3V+\cdots
  }
  
  Factoring out $V$ from each term gives us the equivalent capacitance for
  $N$ capacitors connected in parallel:

  \eq{-.2in}{
    \boxed{C_\text{parallel}=\sum_{i=1}^N C_i}
  }
\end{frame}



\begin{frame}{Capacitors in Series}
  Likewise, we can do a similar analysis to capacitors connected in series.
  \begin{center}
    \begin{tikzpicture}[scale=1.2,thick]
      \draw (0,0) to[C=$C_1$,o-] (1.25,0) to[C=$C_2$] (2.5,0)
      to[C=$C_3$,-o] (3.75,0);
    \end{tikzpicture}
  \end{center}
  The total voltage across these capacitors are the sum of the voltages across
  each of them, i.e.\ $V_\text{tot}=V_1+V_2+V_3+\cdots$
  
  \vspace{.1in}The charge stored on all the capacitors must be the same! The
  total voltage in terms of capacitance and charge is:

  \eq{-.1in}{
    V_\text{tot}=\frac Q{C_1}+\frac Q{C_2}+\frac Q{C_3}+\cdots
  }
\end{frame}




\begin{frame}{Equivalent Capacitance in Series}
  The inverse of the equivalent capacitance for $N$ capacitors connected in
  series is the sum of the inverses of the individual capacitance.

  \eq{-.1in}{
    \boxed{
      \frac1{C_\text{series}}=\sum_{i=1}^N\frac1{C_i}
    }
  }
  
  Make sure we don't confuse ourselves with resistors\footnote{You should know
  about parallel and series resistors in either Grade 9 Science, AP Physics 2
  or Grade 11 Physics.}
\end{frame}



\begin{frame}{How Do We Know That Charges Are The Same?}
  It's simple to show that the charges across all the capacitors are the same
  \begin{center}
    \begin{tikzpicture}[scale=1.5]
      \draw[thick] (0,0) to[C=$C_1$,o-] (1.25,0) to[C=$C_2$,-o] (2.5,0);
      \draw[dashed] (0.625,-.75) rectangle (1.875,1);
    \end{tikzpicture}
  \end{center}
  The capacitor plates and the wire connecting them are really one piece of
  conductor. There is nowhere for the charges to leave the conductor, therefore
  when charges are accumulating on $C_1$, $C_2$ must also have the same charge
  because of conservation of charges.
\end{frame}

\end{document}
